\section{A finite Euler--Maclaurin construction on the whole plane}
\label{em:section}

The purpose of this section is to construct the continuation, rather than to
invoke an abstract analytic-continuation theorem.  Its inputs are the
primitive recursive approximators for rational arithmetic, logarithm,
complex exponential, and finite integrals developed earlier.  Every
improper integral below comes with a tail bound; it is notation for the
resulting approximation procedure.  Likewise, an equality of named complex
numbers means the family of finite rational error statements at every
requested precision.  No quantification over all functions or an evaluator
for arbitrary primitive recursive program codes is needed.

\subsection{Bernoulli polynomials as rational finite data}

Put $B_0(X)=1$.  For $n\geq1$, define $B_n(X)$ to be the unique rational
polynomial satisfying
\begin{equation}
 B_n'(X)=nB_{n-1}(X),\qquad \int_0^1 B_n(x)\,dx=0.
 \label{em:bernoulli-definition}
\end{equation}
Here integration of a polynomial means its finite rational antiderivative.
Thus this is a primitive recursion on coefficient lists: integrate
$nB_{n-1}$ with zero constant, compute its rational integral on $[0,1]$,
and subtract that integral.  Write $B_n=B_n(0)$.  In particular
\[
 B_0=1,\quad B_1=-\tfrac12,\quad B_2=\tfrac16,\quad B_3=0,
 \quad B_4=-\tfrac1{30},\quad B_6=\tfrac1{42}.
\]
Our convention for $B_1$ matters at the value $\zeta(0)$.

We record the required algebra, including its proofs.  Integrating the
first identity in \eqref{em:bernoulli-definition} gives
\begin{equation}
 B_n(1)=B_n(0)\quad(n\geq2),\qquad B_1(1)-B_1(0)=1.
 \label{em:bernoulli-endpoints}
\end{equation}
The polynomial $(-1)^n B_n(1-X)$ has derivative
$n(-1)^{n-1}B_{n-1}(1-X)$ and integral zero.  Starting with $B_0=1$,
uniqueness in the defining recursion therefore proves
\begin{equation}
 B_n(1-X)=(-1)^nB_n(X).
 \label{em:bernoulli-reflection}
\end{equation}
For odd $n\geq3$, the endpoint and reflection identities give
$B_n=-B_n$, and rational arithmetic gives $B_n=0$.
Repeated differentiation at zero gives the coefficient formula
\begin{equation}
 B_n(X)=\sum_{j=0}^n\binom nj B_jX^{n-j}.
 \label{em:bernoulli-coefficients}
\end{equation}
Finally, differentiation and \eqref{em:bernoulli-endpoints} prove
\begin{equation}
 B_n(X+1)-B_n(X)=nX^{n-1}\qquad(n\geq1).
 \label{em:bernoulli-difference}
\end{equation}
Indeed the derivatives agree by the corresponding identity for $n-1$;
for $n\geq2$ the constants agree by evaluating at zero, while $n=1$
is immediate.  Summing this polynomial identity proves Faulhaber's formula
without any limiting operation:
\begin{equation}
 \sum_{n=1}^{N-1}n^q
   =\frac{B_{q+1}(N)-B_{q+1}(1)}{q+1}
       \qquad(q\geq0,\ N\geq1).
 \label{em:faulhaber}
\end{equation}
All these arguments can instead be run on coefficient lists.  The
invariants are finite rational equalities, so that the inductions used
here have quantifier-free arithmetical instances after coding.

For subsequent bounds define the nonnegative rational
\begin{equation}
 C_n=\sum_{j=0}^n\binom nj|B_j|.
 \label{em:bernoulli-bound}
\end{equation}
Then $|B_n(x)|\leq C_n$ for $0\leq x\leq1$.  This deliberately crude
bound is completely explicit and primitive recursive.  For $n=2$ we
shall also use the sharper bound
\begin{equation}
 B_2(x)=x^2-x+\tfrac16=(x-\tfrac12)^2-\tfrac1{12},
 \qquad |B_2(x)|\leq\tfrac16\quad(0\leq x\leq1).
 \label{em:B2-bound}
\end{equation}

\subsection{The finite summation identity, with its sign checked}

On an interval $[n,n+1]$, write $\widetilde B_j(x)=B_j(x-n)$.
For $j\geq2$ the values from adjacent intervals agree at integers.
For $j=1$, the endpoint convention will not affect an integral.
Let $N<M$ be positive integers, $m\geq1$, and let $f$ have the explicitly
controlled derivatives needed in the following finite integrations.
Integration by parts on one unit interval gives
\[
 \int_n^{n+1}B_1(x-n)f'(x)\,dx
 =\frac{f(n)+f(n+1)}2-\int_n^{n+1}f(x)\,dx.
\]
Sum over $n=N,\ldots,M-1$.  Then repeatedly use
$B_{j+1}'=(j+1)B_j$.  The odd Bernoulli endpoint values beyond $B_1$
vanish, while every even endpoint value equals $B_{2k}$.  This gives the
finite identity
\begin{align}
 \sum_{n=N}^{M-1}f(n)
 ={}&\int_N^M f(x)\,dx+\frac{f(N)-f(M)}2 \notag\\
 &+\sum_{k=1}^m\frac{B_{2k}}{(2k)!}
          \bigl(f^{(2k-1)}(M)-f^{(2k-1)}(N)\bigr)\notag\\
 &-\frac1{(2m)!}\int_N^M\widetilde B_{2m}(x)f^{(2m)}(x)\,dx.
 \label{em:finite}
\end{align}
There is no infinite summation in this assertion.  The number of
integrations by parts is $2m-1$; the number of unit intervals is $M-N$.
Both are finite loops.  If integrals are replaced by rational quadrature,
the identity means that its defect is bounded by the specified quadrature
errors; making those errors smaller than $2^{-r}$ proves its named-number
version.  It does not require deciding whether a real error is zero.

Write $(s)_j=\prod_{a=0}^{j-1}(s+a)$, with $(s)_0=1$.  For
$f(x)=x^{-s}=\exp(-s\log x)$,
\[
 f^{(j)}(x)=(-1)^j(s)_j x^{-s-j}.
\]
For $s$ supplied with an apartness bound from $1$, define
\begin{align}
 F_{m,N}(s)={}&\sum_{n=1}^{N-1}n^{-s}
 +\frac{N^{1-s}}{s-1}+\frac12N^{-s}\notag\\
 &+\sum_{k=1}^m\frac{B_{2k}}{(2k)!}
                    (s)_{2k-1}N^{1-s-2k}.
 \label{em:F}
\end{align}
Substitution into \eqref{em:finite} gives exactly
\begin{equation}
 F_{m,N}(s)-F_{m,M}(s)
  =\frac{(s)_{2m}}{(2m)!}
       \int_N^M\widetilde B_{2m}(x)x^{-s-2m}\,dx.
 \label{em:difference}
\end{equation}
In particular, the remainder in the infinite-tail formula below has a
minus sign.

\subsection{Removing the apparent division at $s=1$}

Define an entire elementary function, by its power series or its finite
integral with an explicit exponential modulus,
\begin{equation}
 Q_N(s)=-(\log N)\int_0^1
                \exp\bigl(-t(s-1)\log N\bigr)\,dt.
 \label{em:Q}
\end{equation}
The fundamental theorem of calculus, or termwise finite polynomial
integration followed by the exponential tail estimate, gives
\[
 Q_N(s)=\frac{N^{1-s}-1}{s-1}\quad(s\mathrel{\#}1),
 \qquad Q_N(1)=-\log N.
\]
Here $s\mathrel{\#}1$ means an explicitly supplied positive rational
lower bound for $|s-1|$; it is stronger data than a bare negated equality.
Set
\begin{align}
 E_{m,N}(s)={}&\sum_{n=1}^{N-1}n^{-s}+Q_N(s)+\frac12N^{-s}\notag\\
 &+\sum_{k=1}^m\frac{B_{2k}}{(2k)!}
                  (s)_{2k-1}N^{1-s-2k}.
 \label{em:E}
\end{align}
This expression is defined even at $s=1$.  The finite calculation leading
to \eqref{em:difference} can be performed with $Q_N-Q_M$ in place of
the quotient.  Consequently
\begin{equation}
 E_{m,N}(s)-E_{m,M}(s)
  =\frac{(s)_{2m}}{(2m)!}
     \int_N^M\widetilde B_{2m}(x)x^{-s-2m}\,dx
 \label{em:E-difference}
\end{equation}
holds for every $s$, including $s=1$, without a continuity or identity
theorem argument.

Let $\sigma=\operatorname{Re}s$.  If $\sigma+2m-1>0$, absolute values
in \eqref{em:E-difference} give
\begin{equation}
 |E_{m,N}(s)-E_{m,M}(s)|
 \leq \frac{|(s)_{2m}|C_{2m}}{(2m)!}
       \frac{N^{1-\sigma-2m}}{\sigma+2m-1}.
 \label{em:cauchy-bound}
\end{equation}
The estimate follows by integrating $x^{-\sigma-2m}$; using a positive
rational lower bound for the denominator gives a rational error bound.
Thus $E_{m,N}$ has an explicitly named limit $H(s)$ on this half-plane.
For any such $m$ and any $N\geq1$,
\begin{align}
 H(s)={}&E_{m,N}(s)
 -\frac{(s)_{2m}}{(2m)!}
       \int_N^\infty\widetilde B_{2m}(x)x^{-s-2m}\,dx.
 \label{em:H-integral}
\end{align}
This is also a definition of the improper integral: truncate at $M$ and
use \eqref{em:cauchy-bound} with $N=M$ as its tail estimate.

Here is one explicit, uniform primitive recursive modulus.  Suppose the
input includes integers $A\geq0$ and $R\geq1$ with
\[
 \operatorname{Re}s\geq-A,\qquad |s|\leq R.
\]
Put $m=A+2$, $L=2m$, $U=R+L$, and
\begin{equation}
 K_0=\frac{C_L U^L}{L!},\qquad
 N(k)=2^{k+2}\bigl(1+\lceil K_0\rceil\bigr).
 \label{em:PR-modulus}
\end{equation}
Since $L-A-1\geq1$, \eqref{em:cauchy-bound} is at most $K_0/N$.
It follows that
\begin{equation}
 |H(s)-E_{m,N(k)}(s)|\leq2^{-k-2}.
 \label{em:explicit-evaluation}
\end{equation}
Evaluate the finite expression with rational error at most $2^{-k-2}$;
its result approximates $H(s)$ within $2^{-k-1}$.  The precision budget
for its subexpressions is obtained from finite absolute-value and
Lipschitz bounds for the displayed arithmetic expression.  It consists
only of a bounded sum, a bounded product, and evaluations of the fixed
logarithm and exponential families.  Every loop bound and every requested
subprecision is primitive recursive in $A,R,k$ and the supplied input
name parameters.

There is no dependence on the chosen chart.  For adjacent orders,
\[
 E_{m+1,N}(s)-E_{m,N}(s)
 =\frac{B_{2m+2}}{(2m+2)!}(s)_{2m+1}N^{-s-2m-1}.
\]
On the domain $\sigma+2m-1>0$, this tends to zero with an explicit power
bound.  A finite telescoping sum compares any two admissible orders.
Thus the chart values agree to every requested precision.  This is a
finite compatibility proof, not an application of the identity theorem.

For $s\mathrel{\#}1$, define
\begin{equation}
 \zeta(s)=\frac1{s-1}+H(s).
 \label{em:zeta-definition}
\end{equation}
If $\sigma>1$ with a rational margin, the elementary integral estimate
for $\sum n^{-s}$ and the vanishing of every correction in
\eqref{em:F} show directly that this agrees with its Dirichlet series.
Indeed $N^{1-s}/(s-1)$, $N^{-s}$, and every displayed Bernoulli correction
have explicit bounds tending to zero.  Equation \eqref{em:H-integral}
therefore yields the usual Euler--Maclaurin representation of the
Riemann zeta function, with a concrete primitive recursive evaluation
procedure on every bounded chart away from the pole.  Compare the
classical formula in \cite{DLMF25_2}; the argument above proves the
finite estimate needed here.

\subsection{Entire regular part and controlled derivatives}

A useful specialization of \eqref{em:H-integral} is $N=1$:
\begin{equation}
 H(s)=\frac12+
   \sum_{k=1}^m\frac{B_{2k}}{(2k)!}(s)_{2k-1}
  -\frac{(s)_L}{L!}
        \int_1^\infty\widetilde B_L(x)x^{-s-L}\,dx,
 \qquad L=2m.
 \label{em:H-chart}
\end{equation}
We prove analyticity by an explicit local series.  Fix a center $s_0$
with $\operatorname{Re}s_0\geq-A$, and a disk $|w|\leq T$, where $T$
is a positive integer.  Increase the even integer $L$ until
\[
 a=L-A-1\geq2T+1.
\]
Expand $x^{-w}=\exp(-w\log x)$.  The coefficient of $w^r$ in the
integral in \eqref{em:H-chart}, at $s=s_0+w$, is
\begin{equation}
 j_r=\frac{(-1)^r}{r!}
       \int_1^\infty\widetilde B_L(x)
                     x^{-s_0-L}(\log x)^r\,dx.
 \label{em:local-coefficients}
\end{equation}
Finite integration by parts, repeated $r$ times, proves
\begin{equation}
 \int_1^\infty x^{-a-1}(\log x)^r\,dx
       =\frac{r!}{a^{r+1}}\qquad(a>0).
 \label{em:log-moment}
\end{equation}
The boundary terms vanish by the elementary exponential bound
$t^r\leq r!b^{-r}e^{bt}$ with, for example, $b=a/2$.
Consequently
\begin{equation}
 |j_r|\leq\frac{C_L}{a^{r+1}},\qquad
 \sum_{r>d}|j_rw^r|
 \leq\frac{C_L}{a}\,
          \frac{(T/a)^{d+1}}{1-T/a}.
 \label{em:local-tail}
\end{equation}
The equality between the integral and this series follows first for
finite sums by linearity.  The absolute remainder of the exponential
series has an integrable majorant, and summing its geometric integral
bound gives exactly \eqref{em:local-tail}; no unquantified interchange
of two limits is being used.  Multiplication by the finite polynomial
$(s_0+w)_L$ now gives a power series for $H(s_0+w)$ with an explicit
geometric tail on the chosen disk.  Since every center and radius have
such a chart, $H$ is entire in this constructive, local-power-series
sense.

For later use we also give a uniform bound for differentiated
Euler--Maclaurin remainders.  Assume $\sigma\geq-A$, $|s|\leq R$,
$L=2m\geq A+3$, and $U=R+L$.  For $j\leq L$, differentiating a product
shows
\[
 |((s)_L)^{(j)}|\leq L^jU^L;
\]
for $j>L$ the derivative is zero.  For $x\geq1$ and $v\geq0$,
$(\log x)^v\leq v!x$.  Applying these two inequalities to the finite
Leibniz expansion, and then integrating, gives
\begin{equation}
 \left|H^{(d)}(s)-E_{m,N}^{(d)}(s)\right|
 \leq\frac{K_d}{N},\qquad
 K_d=\frac{C_L U^L}{L!}\,d!(1+L)^d.
 \label{em:derivative-tail}
\end{equation}
More explicitly, the integrand after $d$ derivatives is bounded by
\[
 \frac{C_LU^L}{L!}
 \sum_{j=0}^{\min(d,L)}\binom dj L^j(d-j)!\,x^{A-L+1},
\]
and the coefficient sum is at most $d!(1+L)^d$, while
$\int_N^\infty x^{A-L+1}dx\leq1/N$.
Differentiation of the named integrals is justified by the local series
just constructed; it is not an extra assumption about improper
integrals.  This same bound proves convergence of every fixed-order
derivative by a primitive recursive modulus.  The integer $d$ can also
be an input to the explicitly written coefficient and Leibniz loops.

It follows that the only pole of $\zeta$ is at $1$, and that
\[
 (s-1)\zeta(s)=1+(s-1)H(s)
\]
extends to an entire function with value $1$ at $s=1$.
This last statement refers to the displayed regular extension, rather
than an undefined product at the pole.

\subsection{The residue, Euler's constant, and a small zero-free disk}

The pole and its residue have particularly simple quantitative versions.
Use $m=1$ in \eqref{em:H-chart}.  For $|s-1|\leq\tfrac12$,
$\sigma\geq\tfrac12$, $|s|\leq\tfrac32$, and $|s+1|\leq\tfrac52$.
By \eqref{em:B2-bound},
\begin{align*}
 |H(s)|
 &\leq\frac12+\frac{|s|}{12}
       +\frac{|s(s+1)|}{2}
                  \frac{1/6}{\sigma+1}\\
 &\leq\frac12+\frac18+\frac5{24}=\frac56.
\end{align*}
Thus
\begin{equation}
 \left|(s-1)\zeta(s)-1\right|
       \leq\frac56|s-1|\qquad
       (s\mathrel{\#}1,\ |s-1|\leq\tfrac12).
 \label{em:residue-modulus}
\end{equation}
This proves residue $1$ with an explicit linear modulus.  It also gives
\begin{equation}
 |(s-1)\zeta(s)|\geq\frac7{12},
 \qquad |\zeta(s)|\geq\frac7{12|s-1|}
 \label{em:pole-zero-free}
\end{equation}
on that punctured disk.  These are constructive apartness statements.

Let $H_N=\sum_{n=1}^N1/n$.  The logarithmic integral inequalities give,
for $M\geq N\geq1$,
\[
 0\leq(H_N-\log N)-(H_M-\log M)\leq\frac1N.
\]
To see the sign, sum
$1/(n+1)\leq\int_n^{n+1}dx/x\leq1/n$; for the upper bound sum the
differences $1/n-1/(n+1)$.  Hence
\begin{equation}
 \gamma=\lim_{N\to\infty}(H_N-\log N),\qquad
 |\gamma-(H_N-\log N)|\leq\frac1N
 \label{em:gamma-name}
\end{equation}
defines a primitive recursive real.  Direct evaluation of
\eqref{em:E} at $s=1$ gives
\begin{equation}
 E_{m,N}(1)=H_N-\log N-\frac1{2N}
              +\sum_{k=1}^m\frac{B_{2k}}{2k}N^{-2k}.
 \label{em:gamma-identification}
\end{equation}
Every extra term tends to zero with a displayed power bound.  Therefore
$H(1)=\gamma$ and
\begin{equation}
 \zeta(s)=\frac1{s-1}+\gamma+O(s-1)
 \quad\hbox{near }s=1.
 \label{em:laurent-first}
\end{equation}
The $O$ notation here abbreviates the local power-series modulus above;
it is not an unevaluated existence assertion.

\subsection{All nonpositive integer values}

Fix $q\geq0$ and choose $L=2m>q+1$.  At $s=-q$, $(s)_L=0$ exactly,
so the remainder in \eqref{em:H-integral} vanishes exactly.  In
\eqref{em:F}, the surviving correction terms satisfy
\[
 (-q)_{2k-1}=-\frac{q!}{(q-2k+1)!}
          \quad(2k\leq q+1),
\]
and all later terms are zero.  By
\eqref{em:bernoulli-coefficients},
\[
 \frac{B_{q+1}(N)}{q+1}
 =\frac{N^{q+1}}{q+1}-\frac12N^q+
   \sum_{2k\leq q+1}\frac{B_{2k}}{q+1}
                    \binom{q+1}{2k}N^{q+1-2k}.
\]
Substitute Faulhaber's formula \eqref{em:faulhaber} into
\eqref{em:F}.  Every power of $N$ cancels, including the constant term
when $q+1$ is even.  The exact answer is
\begin{equation}
 \zeta(-q)=-\frac{B_{q+1}(1)}{q+1}\qquad(q\geq0).
 \label{em:negative-unified}
\end{equation}
Consequently,
\begin{equation}
 \zeta(0)=-\frac12,\qquad
 \zeta(-q)=-\frac{B_{q+1}}{q+1}\quad(q\geq1).
 \label{em:negative}
\end{equation}
In particular, for $n\geq1$,
\begin{equation}
 \zeta(-2n)=0,
 \label{em:trivial-zeros}
\end{equation}
because $B_{2n+1}=0$.  Some exact examples are
\[
 \zeta(-1)=-\frac1{12},\quad
 \zeta(-3)=\frac1{120},\quad
 \zeta(-5)=-\frac1{252},\quad
 \zeta(-7)=\frac1{240},\quad
 \zeta(-11)=\frac{691}{32760}.
\]
These values agree with \cite{DLMF25_6}.  Their proof here uses only
finite Bernoulli identities and an exactly vanishing remainder, so the
special-value step is especially direct to arithmetize.  It makes no
claim that the negative even zeros are simple; simplicity requires
additional information, supplied separately if a functional equation
and the corresponding nonvanishing factors have been established.

\subsection{Every Stieltjes coefficient and a global Laurent expansion}

For $r\geq0$ and $N\geq1$, put
\begin{equation}
 c_{r,N}=\sum_{n=1}^N\frac{(\log n)^r}{n}
           -\frac{(\log N)^{r+1}}{r+1}.
 \label{em:stieltjes-approximant}
\end{equation}
For $r=0$ the convention $(\log n)^0=1$ is used, including at $n=1$.
We give a convergence proof with a modulus independent of any compactness
principle.  Set $f_r(x)=(\log x)^r/x$.  For each integer $n\geq1$,
\[
 \left|f_r(n)-\int_n^{n+1}f_r(x)\,dx\right|
       \leq\int_n^{n+1}|f_r'(x)|\,dx.
\]
The inequality follows by integrating
$f_r(n)-f_r(x)=-\int_n^x f_r'(t)dt$ and using the unit interval length.
For $r\geq1$,
\[
 |f_r'(x)|\leq
 \frac{r(\log x)^{r-1}+(\log x)^r}{x^2};
\]
for $r=0$ use $|f_0'(x)|=x^{-2}$.
Let
\[
 I_r(N)=\int_N^\infty\frac{(\log x)^r}{x^2}dx
       =\frac1N\sum_{j=0}^r\frac{r!}{(r-j)!}(\log N)^{r-j}.
\]
The equality is repeated integration by parts with its exponential
boundary bound.  Moreover
$I_r(N)=f_r(N)+rI_{r-1}(N)$ when $r\geq1$.
The absolutely controlled sum of the cell defects defines a real
$\gamma_r$.  Since $f_r(N)\to0$, comparison with
\eqref{em:stieltjes-approximant} gives
\begin{equation}
 |\gamma_r-c_{r,N}|\leq2I_r(N)
 \leq\frac{2^{r+2}r!}{\sqrt N}.
 \label{em:stieltjes-tail}
\end{equation}
For the last inequality use
$(\log x)^r\leq2^r r!\sqrt x$ and integrate $x^{-3/2}$.
Thus, for example,
\begin{equation}
 N_r(k)=2^{2(k+r+3)}(r!)^2
 \label{em:stieltjes-modulus}
\end{equation}
makes the truncation error at most $2^{-k-1}$.  The formula is primitive
recursive in both $r$ and $k$, and the required rational approximations
to the finitely many logarithms have an explicit finite error budget.
For $r=0$ this is the same $\gamma$ as before.

We next identify these numbers with the Laurent coefficients.  Finite
differentiation of \eqref{em:E} at $s=1$ gives
\begin{align}
 E_{m,N}^{(r)}(1)={}&(-1)^r c_{r,N}
       -\frac{(-1)^r(\log N)^r}{2N}\notag\\
 &+\left.\frac{d^r}{ds^r}
      \sum_{k=1}^m\frac{B_{2k}}{(2k)!}
                     (s)_{2k-1}N^{1-s-2k}\right|_{s=1}.
 \label{em:stieltjes-derivative}
\end{align}
In particular the derivative of $Q_N$ contributes
$(-1)^{r+1}(\log N)^{r+1}/(r+1)$; this follows directly by integrating
$t^r$ in \eqref{em:Q}.  Every term on the second line of
\eqref{em:stieltjes-derivative} is a finite rational linear combination
of $(\log N)^v/N^{2k}$ with $v\leq r$.  The bound
$(\log N)^v\leq2^v v!\sqrt N$ supplies a primitive recursive vanishing
modulus for all of them.  Equation \eqref{em:derivative-tail}, with any
fixed admissible $m$, therefore proves
\begin{equation}
 H^{(r)}(1)=(-1)^r\gamma_r.
 \label{em:stieltjes-identification}
\end{equation}

The global Laurent expansion is consequently
\begin{equation}
 \zeta(1+w)=\frac1w+
        \sum_{r=0}^\infty\frac{(-1)^r\gamma_r}{r!}w^r
        \qquad(w\mathrel{\#}0).
 \label{em:global-laurent}
\end{equation}
This is valid at every supplied bounded value of $w$, with an explicit
geometric modulus for the regular series.  To make that assertion
concrete, fix an integer $T\geq1$ with $|w|\leq T$, and choose an even
$L>2T$.  In \eqref{em:local-coefficients}, take $s_0=1$; then
$|j_r|\leq C_L/L^{r+1}$.  The other factor is the degree-$L$ polynomial
$P(w)=(1+w)_L$.  Its coefficient absolute values at $T$ sum to
$P(T)=\prod_{a=1}^L(T+a)\leq(1+T+L)^L$.
For $h\geq L$, the omitted terms of the product series are therefore
bounded by
\begin{equation}
 \left|H(1+w)-\sum_{r=0}^h
           \frac{(-1)^r\gamma_r}{r!}w^r\right|
 \leq
 \frac{C_L(1+T+L)^L}{L!L}\,
       \frac{(T/L)^{h-L+1}}{1-T/L}.
 \label{em:global-laurent-tail}
\end{equation}
Indeed a polynomial coefficient of degree $\ell\leq L$ multiplies only
integral-series terms of degree greater than $h-\ell$ in the omitted
tail, and this is bounded by the geometric tail starting after $h-L$.
The polynomial part of \eqref{em:H-chart} has degree at most $L-1$ and
contributes no remainder.  Since $T/L<1/2$, this is a primitive recursive
geometric rate.  No assertion that an unspecified analytic power series
has a computable radius is used.

\subsection{Further consequences with explicit bounds}

\paragraph{Conjugation.}
Every Bernoulli coefficient is rational and every positive-base power
satisfies $\overline{x^{-s}}=x^{-\overline s}$.  Apply conjugation first
to the finite expressions \eqref{em:E}, then use their common error
bound.  This proves
\begin{equation}
 H(\overline s)=\overline{H(s)},\qquad
 \zeta(\overline s)=\overline{\zeta(s)}\quad(s\mathrel{\#}1).
 \label{em:conjugation}
\end{equation}
In particular real inputs away from $1$ have real values, and a certified
zero gives its conjugate zero.

\paragraph{An explicit, nonoptimal growth bound.}
Given $\sigma\geq-A$, $|s|\leq R$, $|s-1|\geq2^{-D}$, choose
$L=2m\geq A+2$ and $U=R+L$.  Equation \eqref{em:H-chart} gives
\begin{equation}
 |\zeta(s)|\leq 2^D+\frac12+
  \sum_{k=1}^m\frac{|B_{2k}|}{(2k)!}U^{2k-1}
          +\frac{C_LU^L}{L!}.
 \label{em:growth}
\end{equation}
For each fixed left boundary $-A$, this is a concrete polynomial bound
in $R$ away from the pole.  It is sufficient for approximation and
formalization; it is not the sharp growth estimate studied in analytic
number theory.

\paragraph{Compatibility with the alternating series.}
Let $S_N(s)=\sum_{n=1}^Nn^{-s}$.  Taking $m=1$ in the continuation
formula and moving the $n=N$ term between the sum and the endpoint gives
\begin{align}
 \zeta(s)={}&S_N(s)+\frac{N^{1-s}}{s-1}-\frac12N^{-s}
       +\frac{s}{12}N^{-s-1}\notag\\
 &-\frac{s(s+1)}2\int_N^\infty\widetilde B_2(x)x^{-s-2}dx.
 \label{em:SN}
\end{align}
If $\sigma>0$, $|s|\leq R$, and $s\mathrel{\#}1$, it follows that
\begin{equation}
 \left|S_N(s)-\zeta(s)-\frac{N^{1-s}}{1-s}\right|
 \leq\left(\frac12+\frac R{12}+\frac{R(R+1)}{12}\right)N^{-\sigma}.
 \label{em:eta-compatibility}
\end{equation}
The identity
$\sum_{n=1}^{2N}(-1)^{n-1}n^{-s}
 =S_{2N}(s)-2^{1-s}S_N(s)$ is finite and exact.  Inserting
\eqref{em:eta-compatibility} cancels the two growing powers exactly and
proves the compatibility with the eta representation on its right
half-plane.  Division by $1-2^{1-s}$, when desired, separately requires
an apartness certificate for that denominator.
